\section{Objetivos}

\subsection{Objetivo General}

Analizar y comprender el movimiento de los cuerpos mediante el estudio de la cinemática, 
aplicando modelos matemáticos y herramientas de cálculo para describir, interpretar y resolver 
problemas relacionados con el Movimiento Rectilíneo Uniforme (MRU), el Movimiento Rectilíneo 
Uniformemente Acelerado (MRUA) y la caída libre (MVCL), así como desarrollar la simulación de un 
caso de MRUA mediante el uso del software GeoGebra.

\subsection{Objetivos Específicos}

\begin{itemize}
    \item Describir las magnitudes fundamentales de la cinemática, tales como posición, desplazamiento, velocidad y aceleración.
    \item Analizar las características del Movimiento Rectilíneo Uniforme (MRU) y del Movimiento Rectilíneo Uniformemente Acelerado (MRUA).
    \item Aplicar conceptos de cálculo diferencial e integral en la obtención de ecuaciones de movimiento.
    \item Estudiar el fenómeno de la caída libre considerando la aceleración de la gravedad.
    \item Resolver problemas físicos mediante el uso de ecuaciones cinemáticas.
    \item Utilizar herramientas tecnológicas para la simulación y verificación de los resultados obtenidos.
\end{itemize}